\chapter*{摘\quad 要}

随着生成式人工智能技术的快速发展，AI生成内容在图像生成、视频生成、数字人驱动与智能创作等领域得到广泛应用，但同时也带来了虚假合成内容传播、不安全内容扩散、身份伪造和平台审核压力增加等问题。围绕AIGC内容在真实应用场景中的安全风险，本文以图像与视频两类典型媒体为研究对象，研究虚假内容检测与不安全内容识别方法，并在此基础上设计实现一个多模态内容安全检测与分析平台。

本文首先梳理AIGC生成技术、深度伪造技术以及安全检测领域的发展现状，分析现有方法在检测对象、泛化能力、部署成本和工程适配性等方面的特点。随后，围绕图像与视频安全分析任务，构建由前端展示层、后端代理层和模型推理层组成的系统架构，整合本地部署检测模型与外部多模态审核接口，形成统一的内容输入、检测分析、风险分级与结果可视化流程。系统支持图像与视频文件上传、URL输入、虚假内容鉴别、不安全内容审核、检测结果展示与内容管理等功能，可为实验评估、样本构建和应用验证提供平台支撑。

在实验与分析部分，本文结合当前系统实现情况，对图像虚假检测、视频内容分析、不安全内容审核和平台可用性进行验证，并总结本地模型与外部接口在不同任务中的适用特点。本文工作兼顾工程实现与应用价值，可为AIGC内容安全治理、网络平台审核以及多模态安全研究提供参考。

\vspace{1em}
\noindent\textbf{关键词：}AIGC安全；多模态检测；虚假内容检测；不安全内容识别；内容安全平台
